%------------------------------------------------------------------------------%
% \chapter{Revisão}
%------------------------------------------------------------------------------%

%------------------------------------------------------------------------------%
\begin{exercicios}{Polinômios}
    
    \begin{exercicio}
        Defina uma função $f(x)$ que forneça a área da região destacada na figura, lembrando que a área de um retângulo de lados $b$ e $h$ é $bh$.
        \addgraphic{.8}{exercicio-polinomio-area-retangulo}
    \end{exercicio}
    
    \begin{exercicio}
        Encontre as raízes de cada equação por fatoração.
        \begin{mathlist}[2]
            \item x^2 + x - 12 = 0
            \item 3x^2 - x - 4 = 0
            \item \frac{1}{4} x^2 - x + 1 = 0
            \item 8x^2 + 2x - 3 = 0
            \item -6x^2 - 10x + 4 = 0
            \item 2x^4 + x^2 = 1
        \end{mathlist}
        \begin{answers}
            \begin{mathlist}[2]
                \item x_1 = -4 \qquad x_2 = 3
                \item x_1 = -1 \qquad x_2 = \sfrac{4}{3}
                \item x_1 = 2
            \end{mathlist}
        \end{answers}
    \end{exercicio}
    
    \begin{exercicio}
        Resolva as equações utilizando a fórmula quadrática.
        \begin{mathlist}[2]
            \item 4x^2 + 5x - 6 = 0
            \item 8x^2 - 8x - 3 = 0
            \item x^2 - 6x + 6 = 0
            \item 2x^2 + 8x + 7 = 0
        \end{mathlist}
        \begin{answers}
            \begin{mathlist}[2]
                \item x_1 = -2 \qquad x_2 = \sfrac{4}{3}
                \item x_{1,2} = \frac{2\pm\sqrt{10}}{4}
                \item x_{1,2} = 3\pm\sqrt{3}
                \item x_{1,2} = -2\pm\frac{\sqrt{2}}{2}
            \end{mathlist}
        \end{answers}    
    \end{exercicio}
    
    \begin{exercicio}
        Encontre as raízes da equação $2x^4 + x^2 = 1$ substituindo $x^2$ por $y$.
        \begin{answers}
            \(\qquad x_{1,2} = \pm\frac{1}{\sqrt{2}} \)
        \end{answers}
    \end{exercicio}    
    
    \begin{exercicio}
        Resolva as inequações quadráticas.
        \begin{mathlist}[2]
            \item   x^2 + 3x \geq 10
            \item -3x^2 - 11x + 4 > 0
            \item -4x^2 + 4x - 1 < 0
            \item   x^2 + x + 2 \leq 0
        \end{mathlist}
        \begin{answers}
            \begin{mathlist}[2]
                \item x \leq -5 \text{ ou } x \geq 2
                \item -4 < x < 1/3
                \item x \neq 1/2
                \item \) Não há solução.
            \end{mathlist}
        \end{answers}
    \end{exercicio}
    
    \begin{exercicio}
        Para cada expressão na forma $p(x)/d(x)$ abaixo, calcule o quociente $q(x)$ e o resto $r(x)$.
        \begin{mathlist}
            \item (2x^3 - 3x^2 + 6)/(x^2 - 2)
            \item (6x^2 - 4x - 3)/(3x - 5)
            \item (x^4 + 2x - 12)/(x + 2)
            \item (4x^3 + 2x^2 + 11x)/(2x^2 + 3)
            \item (6x^4 + 5x^3 - 2x)/(3x - 2)
            \item (4x^3 + 6x - 10)/(2x - 4)
            \item (x^2 - 5x + 8)/(x - 3)
            \item (3x + 7)/(x + 4)
            \item (x^4 - 2)/(x - 1)
            \item (24x^3 - 4x - 1)/(2x - 1)
            \item (8x^3 - 12x^2 - 2x)/(4x - 8)
            \item (x^3 - 3x^2 + 4x - 5)/(x - 4)
            \item (2x^4 - 4x^3 + x - 17)/(x^2 - 4)
            \item (x^4 - 6x^3 + 3x^2 - 2x + 3)/(x^2 - 2x - 3)
            \item (x^4 - 5x^2 + 4)/(x^2 - 1)
            \item (3x^5 - 2x^3 - 11x)/(x^3 - 3x)
            \item (6x^2 + 7x + 9)/(2x^2 - 5x + 1)
        \end{mathlist}
        \begin{answers}
            \begin{alphalist}
                \item $q(x) = 2x - 3                $, $r(x) = 4x$
                \item $q(x) = 2x + 2                $, $r(x) = 7$
                \item $q(x) = x^3 - 2x^2 + 4x - 6   $, $r(x) = 0$
                \item $q(x) = 2x + 1                $, $r(x) = 5x - 3$
                \item $q(x) = 2x^3 + 3x^2 + 2x + 2/3$, $r(x) = 4/3$
                \item $q(x) = 2x^2 + 4x + 11        $, $r(x) = 34$
                \item $q(x) = x - 2                 $, $r(x) = 2$
                \item $q(x) = 3                     $, $r(x) = -5$
                \item $q(x) = x^3 + x^2 + x + 1     $, $r(x) = -1$
                \item $q(x) = 12x^2 + 6x + 1        $, $r(x) = 0$
                \item $q(x) = 2x^2 + x + 3/2        $, $r(x) = 12$
                \item $q(x) = x^2 + x + 8           $, $r(x) = 27$
                \item $q(x) = 2x^2 - 4x + 8         $, $r(x) = 15 - 15x$
                \item $q(x) = x^2 - 4x - 2          $, $r(x) = -18x - 3$
                \item $q(x) = x^2 - 4               $, $r(x) = 0$
                \item $q(x) = 3x^2 + 7              $, $r(x) = 10x$
                \item $q(x) = 3                     $, $r(x) = 22x + 6$
            \end{alphalist}
        \end{answers}
    \end{exercicio}
    
    \begin{exercicio}
        Para os problemas do exercício anterior, expresse $p(x)$ na forma $d(x)q(x) + r(x)$.
        \begin{answers}
            \begin{mathlist}
                \item p(x) = (x^2 - 2)(2x - 3) + 4x
                \item p(x) = (3x - 5)(2x + 2) + 7
                \item p(x) = (x + 2)(x^3 - 2x^2 + 4x - 6)
                \item p(x) = (2x^2 + 3)(2x + 1) + 5x - 3
                \item p(x) = (3x-2)(2x^3 +3x^2 +2x+2/3)+4/3
                \item p(x) = (2x - 4)(2x^2 + 4x + 11) + 34
                \item p(x) = (x - 3)(x - 2) + 2
                \item p(x) = (x + 4)3 - 5
                \item p(x) = (x - 1)(x^3 + x^2 + x + 1) - 1
                \item p(x) = (2x - 1)(12x^2 + 6x + 1)
                \item p(x) = (4x - 8)(2x^2 + x + 3/2 ) + 12
                \item p(x) = (x - 4)(x^2 + x + 8) + 27
                \item p(x) = (x^2 - 4)(2x^2 - 4x + 8) - 15x + 15
                \item p(x) = (x^2 -2x-3)(x^2 -4x-2)-18x-3
                \item p(x) = (x^2 - 1)(x^2 - 4)
                \item p(x) = (x^3 - 3x)(3x^2 + 7) + 10x
                \item p(x) = (2x^2 - 5x + 1)3 + 22x + 6
            \end{mathlist}
        \end{answers}
    \end{exercicio}
    
    \begin{exercicio}
        Para cada expressão na forma $p(x)/d(x)$ abaixo, calcule o quociente $q(x)$ e o resto $r(x)$ usando o algoritmo de Ruffini.
        \begin{mathlist}
            \item (x^4 + 2x - 12)/(x + 2)
            \item (3x^2 + 2x - 5)/(x - 2)
            \item (4x^4 + 6x^3 - 8x^2 + 22x - 24)/(x + 3)
            \item (-2x^3 + 3x^2 + 12x + 25)/(x - 4)
            \item (x^5 - 9x^3 + 2x)/(x - 3)
            \item (-6x^3 + 4x^2 - x + 2)/(x - 1/3)
            \item (2x^3 - 9x^2 + 6x + 5)/(x - 3/2)
            \item (x^2 - 5x - 6)/(x + 1)
            \item (-4x^2 + 11x + 26)/(x - 4)
            \item (6x^2 - 7x - 9)/(x + 1/2)
            \item (x^3 - 9x^2 )/(x - 3)
            \item (5x^4 - 1)/(x - 2)
            \item (8x^4 + 6x^3 + 3x^2 + 1)/(x - 1/2)
            \item (x^4 - 20x^2 - 50)/(x - 5)
            \item x^4 /(x - 3)
            \item (2x^5 - 4x^4 + 9x^3 - 5x^2 + x - 3)/(x - 1)
        \end{mathlist}
        \begin{answers}
            \begin{alphalist}
                \item $q(x) = x^3 - 2x^2 + 4x - 6$,        $r(x) = 0$
                \item $q(x) = 3x + 8$,                     $r(x) = 11$
                \item $q(x) = 4x^3 - 6x^2 + 10x - 8$,      $r(x) = 0$
                \item $q(x) = -2x^2 - 5x - 8$,             $r(x) = -7$
                \item $q(x) = x^4 + 3x^3 + 2$,             $r(x) = 6$
                \item $q(x) = -6x^2 + 2x - 1/3$,           $r(x) = 17/9$
                \item $q(x) = 2x^2 - 6x - 3$,              $r(x) = 1/2$
                \item $q(x) = x - 6$,                      $r(x) = 0$
                \item $q(x) = -4x - 5$,                    $r(x) = 6$
                \item $q(x) = 6x - 10$,                    $r(x) = -4$
                \item $q(x) = x^2 - 6x - 18$,              $r(x) = -54$
                \item $q(x) = 5x^3 + 10x^2 + 20x + 40$,    $r(x) = 79$
                \item $q(x) = 8x^3 + 10x^2 + 8x + 4$,      $r(x) = 3$
                \item $q(x) = x^3 + 5x^2 + 5x + 25$,       $r(x) = 75$
                \item $q(x) = x^3 + 3x^2 + 9x + 27$,       $r(x) = 81$
                \item $q(x) = 2x^4 -2x^3 +7x^2 +2x+3$,     $r(x) = 0$
            \end{alphalist}
        \end{answers}
    \end{exercicio}
    
    \begin{exercicio}
        Para os problemas do exercício anterior, expresse $p(x)/d(x)$ na forma $q(x) + r(x)/d(x)$.
        \begin{answers}
            \begin{mathlist}
                \item p(x)/d(x) = x^3 - 2x^2 + 4x - 6
                \item p(x)/d(x) = 3x + 8 + 11/(x - 2)
                \item p(x)/d(x) = 4x^3 - 6x^2 + 10x - 8
                \item p(x)/d(x) = -2x^2 - 5x - 8 - 7/(x - 4)
                \item p(x)/d(x) = x^4 + 3x^3 + 2 + 6/(x - 3)
                \item p(x)/d(x) = -6x^2 + 2x -1/3+ 17/(9x - 3)
                \item p(x)/d(x) = 2x^2 - 6x - 3 + 1/(2x - 3)
                \item p(x)/d(x) = x - 6
                \item p(x)/d(x) = -4x - 5 + 6/(x - 4)
                \item p(x)/d(x) = 6x - 10 - 4/(x + 12 )
                \item p(x)/d(x) = x^2 - 6x - 18 - 54/(x - 3)
                \item p(x)/d(x) = 5x^3 + 10x^2 + 20x + 40 + 79/(x - 2)
                \item p(x)/d(x) = 8x^3 +10x^2 +8x+4+3/(x- 1/2 )
                \item p(x)/d(x) = x^3 +5x^2 +5x+25+75/(x-5)
                \item p(x)/d(x) = x^3 +3x^2 +9x+27+81/(x-3)
                \item p(x)/d(x) = 2x^4 - 2x^3 + 7x 2 + 2x + 3
            \end{mathlist}
        \end{answers}
    \end{exercicio}
    
    \begin{exercicio}
        Usando o algoritmo de Ruffini, verifique quais valores abaixo correspondem a zeros das funções associadas.
        \begin{alphalist}
            \item $f(x) = x^2 - 3x + 4$,                         $x_1 = 2$,    $x_2 = -2$
            \item $f(x) = -2x^2 + 3x + 2$,                       $x_1 = -1/2$, $x_2 = -2$
            \item $f(x) = 4x + x^2 $,                            $x_1 = -4$,   $x_2 = 0$
            \item $f(x) = -x^2 - 4$,                             $x_1 = 2$,    $x_2 = -2$
            \item $f(x) = x^3 - 4x^2 + 6x - 9$,                  $x_1 = 3$,    $x_2 = 1$
            \item $f(x) = x^4 - 3x^2 + 2$,                       $x_1 = 5$,    $x_2 = -1$
            \item $f(x) = 2x^4 + x^3 - 25x^2 + 12x$,             $x_1 = -4$,   $x_2 = 3$
            \item $f(x) = x^5 + 2x^4 - 3x^3 + 12x^2 - 28x + 16$, $x_1 = 6$,    $x_2 = 2$
            \item $f(x) = 9x^3 - 15x^2 - 26x + 40$,              $x_1 = 1/2$,  $x_2 = 4/3$
            \item $f(x) = x^3 - 21x - 20$,                       $x_1 = 5$,    $x_2 = 1$
        \end{alphalist}
        \begin{answers}
            \begin{alphalist}
                \item Nenhum valor é um zero da função.
                \item Só -1/2 é um zero de $f$
                \item -4 e 0 são zeros de $f$
                \item Nenhum valor é um zero da função.
                \item Só 3 é um zero de $f$
                \item Só -1 é um zero de $f$
                \item -4 e 3 são zeros de $f$
                \item Nenhum valor é um zero da função.
                \item Só 4/3 é um zero de $f$
                \item Só 5 é um zero de $f$
            \end{alphalist}
        \end{answers}
    \end{exercicio}
       
    \begin{exercicio}
        Para cada função polinomial abaixo, determine o valor da constante $c$ de modo que o termo fornecido seja um fator de $p$.
        \begin{alphalist}
            \item $p(x) = x^2 - 9x + c$,                   \tabto{55mm} fator: $x-8$
            \item $p(x) = 5x^2 + cx + 9$,                  \tabto{55mm} fator: $x+3$
            \item $p(x) = x^3 - 6x^2 + 3x + c$,            \tabto{55mm} fator: $x-5$
            \item $p(x) = 3x^3 + cx^2 - 13x + 3$,          \tabto{55mm} fator: $x-1$
            \item $p(x) = x^4 - 2x^3 + 8x^2 + cx - 2$,     \tabto{55mm} fator: $x-2$
            \item $p(x) = 2x^4 - 10x^3 + cx^2 + 6x + 40$,  \tabto{55mm} fator: $x-4$
        \end{alphalist}
        \begin{answers}
            \begin{mathlist}[3]
                \item c = 8
                \item c = 18
                \item c = 10
                \item c = 7
                \item c = -15
                \item c = 4
            \end{mathlist}
        \end{answers}
    \end{exercicio}
    
    \begin{exercicio}
        Determine as raízes das equações abaixo. Escreva na forma fatorada os polinômios que aparecem no lado esquerdo das equações.
        \begin{mathlist}[2]
            \item   x^3 -  4x           = 0
            \item   x^3 -  4x^2 - 21x   = 0
            \item  2x^3 + 11x^2 -  6x   = 0
            \item -3x^3 +  6x^2 +  9x   = 0
            \item   x^4 -   x^3 - 20x^2 = 0
            \item   x^4 -  8x^3 + 16x^2 = 0
            \item  5x^4 -  8x^3 +  3x^2 = 0
            \item  8x^4 -  6x^3 -  2x^2 = 0
        \end{mathlist}
        \begin{answers}
            \begin{alphalist}
                \item Raízes: 0, 2 e -2,                 \\[1mm] $p(x) = x(x - 2)(x + 2)$
                \item Raízes: 0, -3 e 7,                 \\[1mm] $p(x) = x(x + 3)(x - 7)$
                \item Raízes: 0, 1/2 e -6,               \\[1mm] $p(x) = 2x (x - 1/2 ) (x + 6)$
                \item Raízes: 0, 3 e -1,                 \\[1mm] $p(x) = -3x(x - 3)(x + 1)$
                \item Raízes: 5, -4 e 0 (mult. 2),       \\[1mm] $p(x) =  x^2 (x - 5)(x + 4)$
                \item Raízes: 4 (mult. 2) e 0 (mult. 2), \\[1mm] $p(x) =  x^2 (x - 4)^2$
                \item Raízes: 3/5, 1 e 0 (mult. 2),      \\[1mm] $p(x) = 5x^2 (x - 3/5 ) (x - 1)$
                \item Raízes: 3/2, -1/4 e 0 (mult. 2),   \\[1mm] $p(x) = 8x^2 (x - 3/2 ) (x + 1/4)$
            \end{alphalist}
        \end{answers}
    \end{exercicio}
    
    \begin{exercicio}
        Determine as raízes das equações abaixo. Escreva na forma fatorada os polinômios que aparecem no lado esquerdo das equações.
        \begin{alphalist}
            \item $ x^3 +   x^2 -  2x   - 2         = 0$, \\[1mm] sabendo que $x =  -1$              é uma raiz.
            \item $ x^3 -  5x^2 -  4x   + 20        = 0$, \\[1mm] sabendo que $x =   2$              é uma raiz.
            \item $ x^4 -  9x^3 -   x^2 + 81x - 72  = 0$, \\[1mm] sabendo que $x =   8$ e $x = 3$    são raízes.
            \item $ x^3 -  3x^2 - 10x   + 24        = 0$, \\[1mm] sabendo que $x =   4$              é uma raiz.
            \item $ x^3 -  4x^2 - 17x   + 60        = 0$, \\[1mm] sabendo que $x =   3$              é uma raiz.
            \item $4x^4 - 21x^3 - 19x^2 + 6x        = 0$, \\[1mm] sabendo que $x = 1/4$              é uma raiz.
            \item $4x^3 - 16x^2 + 21x   - 9         = 0$, \\[1mm] sabendo que $x =   1$              é uma raiz.
            \item $3x^3 - 26x^2 + 33x   + 14        = 0$, \\[1mm] sabendo que $x =   7$              é uma raiz.
            \item $ x^4 -  9x^3 + 17x^2 + 33x - 90  = 0$, \\[1mm] sabendo que $x =  -2$ e $x = 5$    são raízes.
            \item $ x^4 -  6x^3 -  5x^2 + 30x       = 0$, \\[1mm] sabendo que $x =   6$              é uma raiz.
            \item $2x^4 +  9x^3 - 80x^2 + 21x + 108 = 0$, \\[1mm] sabendo que $x =   4$ e $x = 3/2$  são raízes.
            \item $ x^3 +  7x^2 + 13x   + 15        = 0$, \\[1mm] sabendo que $x =  -5$              é uma raiz.
            \item $3x^3 +  2x^2 + 17x   - 6         = 0$, \\[1mm] sabendo que $x = 1/3$              é uma raiz.
            \item $ x^3 +  7x^2 + 20x   + 32        = 0$, \\[1mm] sabendo que $x =  -4$              é uma raiz.
            \item $ x^3 -  3x^2 +  9x   - 27        = 0$, \\[1mm] sabendo que $x =   3$              é uma raiz.
            \item $2x^3 - 10x^2 - 13x   - 105       = 0$, \\[1mm] sabendo que $x =   7$              é uma raiz.
            \item $ x^4 +  2x^3 -  5x^2 - 36x + 60  = 0$, \\[1mm] sabendo que $x =   2$              é uma raiz de multiplicidade 2.
            \item $ x^4 -  6x^3 + 25x^2 - 150x      = 0$, \\[1mm] sabendo que $x =   6$              é uma raiz.
            \item $6x^4 +  7x^3 +  6x^2 - 1         = 0$, \\[1mm] sabendo que $x = 1/3$ e $x =-1/2$  são raízes.
        \end{alphalist}
        \begin{answers}
            \begin{alphalist}
                \item $p(x) = (x + 1)(x - \sqrt{2})(x + \sqrt{2})$,   \\[1mm] Raízes: -1, $\sqrt{2}$ e $-\sqrt{2}$        
                \item $p(x) = (x - 1)(x - 2)(x - 8)$,                 \\[1mm] Raízes: 1, 2 e 8
                \item $p(x) = (x + 3)(x - 1)(x - 3)(x - 8)$,          \\[1mm] Raízes: -3, 1, 3 e 8             
                \item $p(x) = (x + 3)(x - 2)(x - 4)$,                 \\[1mm] Raízes: -3, 2 e 4
                \item $p(x) = (x + 4)(x - 3)(x - 5)$,                 \\[1mm] Raízes: -4, 3 e 5
                \item $p(x) = 4(x + 1/4 )x(x + 1)(x - 6)$,            \\[1mm] Raízes: -1, 0, 1/4 e 6
                \item $p(x) = 4 (x - 3/2 ) (x - 1)$,                  \\[1mm] Raízes: 1 e 3/2
                \item $p(x) = (x + 5) (x 2 + 2x + 3)$,                \\[1mm] Raiz: -5
                \item $p(x) = 3 (x - 1/3)(x^2+x+6)$,                  \\[1mm] Raiz: 1/3
                \item $p(x) = (x + 4) (x^2 + 3x + 8)$,                \\[1mm] Raiz: -4
                \item $p(x) = (x - 3) (x^2 + 9)$,                     \\[1mm] Raiz: 3
                \item $p(x) = (x - 7) (2x^2 + 4x + 15)$,              \\[1mm] Raiz: 7
                \item $p(x) = (x - 2)^2 (x^2 + 6x + 15)$,             \\[1mm] Raiz: 2 (multiplicidade 2)
                \item $p(x) = x(x - 6) (x^2 + 25)$,                   \\[1mm] Raizes: 0 e 6
                \item $p(x) = 6 (x + 1/2 ) (x - 1/3 ) (x^2 + x + 1)$, \\[1mm] Raizes: -1/2 e 1/3                            
            \end{alphalist}
        \end{answers}
    \end{exercicio}
    
\end{exercicios}

%------------------------------------------------------------------------------%
